
% 12. The Cosmic Paradox of Fertility Collapse and Gender Polarization: A Fatal Collision of Smoothing Trajectories.tex

\section{저출산과 젠더 양극화의 우주적 역설: 평활화 궤적의 치명적 충돌}
최근 한국을 비롯해 북미, 서유럽, 동아시아 등 전 세계의 고도로 발달한 선진국들에서는 매우 기괴하고도 공통적인 현상이 관찰되고 있다. 젊은 여성들은 갈수록 급격하게 진보적인 정치 성향을 띠는 반면, 젊은 남성들은 상대적으로 보수화되며 서로 극단적으로 멀어지는 '젠더 양극화(Gender Polarization)' 현상이다. 사회학자들과 언론은 이를 단순한 인터넷의 선동, 세대 갈등, 혹은 얄팍한 이념 전쟁의 결과로 취급하며 어떻게든 봉합하려 든다.

그러나 디스턴스(DISTANCE)의 서늘하고 무자비한 열역학적 렌즈를 들이대면, 이 현상의 기저에 깔린 진짜 물리학적 실체가 폭로된다. 이 거대한 양극화는 정치의 문제가 아니다. 그것은 완벽한 '평활화(Equalization)'를 향해 맹렬하게 추락하고 있는 문명의 거대한 궤적과, 결코 깎아낼 수 없는 인간의 '생물학적 번식 구조'가 정면충돌하며 내뿜는 뼈아픈 '물리학적 파열음'이다.

앞선 제11편에서 우리는 '깨끗한 유리창의 얼룩'이라는 비유를 통해 문명의 평활화가 빚어낸 역설을 확인했다. 신분, 계급, 교육, 정치적 억압이라는 수많은 더러운 얼룩(불평등)들이 가득했던 과거의 창문에서는, 특정한 차별 하나가 특별히 부각되지 않았다. 그러나 문명이라는 거대한 평활화 엔진이 그 창문을 역사상 가장 투명하고 매끄럽게 닦아내자, 모든 차별이 사라진 그 깨끗한 무대 위에 홀로 남겨진 단 하나의 선천적 비대칭성, 즉 '임신과 출산'이라는 생물학적 얼룩이 누구에게나 섬뜩할 정도로 거대하고 선명하게 발견되기 시작한 것이다.

오늘날의 남성과 여성은 인류 역사상 전대미문의 완벽한 '대칭적 단일 경쟁 트랙' 위로 던져졌다. 남녀는 같은 교실에서 공부하고, 같은 시험을 치르며, 같은 직장에 지원하여 동일한 승진 체계와 성공의 잣대 안에서 서로의 모멘텀을 맹렬하게 부딪치며 경쟁한다. 사회적 역할의 모든 층위가 완벽하게 대칭화된 것이다. 이것은 분명 문명이 이룩한 가장 숭고한 도덕적 성취이자 위대한 진보다.

그런데 바로 그 완벽한 대칭의 투기장 한가운데서, 결코 깎아낼 수 없는 가장 차가운 물리적 팩트가 앞길을 가로막는다. 인간의 번식은 오직 여성의 몸을 통해서만 이루어지며, 출산 이후의 육체적 회복 역시 온전히 여성의 몸을 통해서만 이루어진다는 절대적 비대칭성이다. 어떠한 혁명적 정치 제도도, 천문학적인 보조금 정책도, 고도의 첨단 기술도 이 생물학적 현실 자체를 완전히 소거하거나 남녀에게 완벽히 5대5로 재분배할 수는 없다.

모든 것이 대칭적인 공정 경쟁의 무대 위에서, 특정 성별의 육체에만 독점적으로 부과되는 이 막대한 모멘텀의 소모(출산)는 더 이상 과거처럼 자연의 섭리나 고결한 희생으로 포장되지 않는다. 그것은 동일한 트랙을 달려야 하는 경쟁 주체에게 가해지는 치명적인 '기회비용'이자, 공정성이라는 현대 사회의 가장 신성한 룰을 정면으로 위반하는 거대한 '선천적 불공정'으로 감각되기 시작한 것이다.

젊은 여성들이 갈수록 진보적인 성향을 띠게 되는 역학적 이유가 바로 여기에 있다. 진보란 본질적으로 '더 많은 불평등의 해소와 평활화'를 요구하는 방향성이다. 가장 고도로 평등해진 사회에서 살아가는 여성들은, 공정성의 혜택을 누구보다 완벽하게 체감하고 있기에 오히려 '끝내 평활화되지 않고 남아 있는 유일한 생물학적 비대칭성(출산)'에 대해 가장 극도로 예민한 감각을 지니게 된다. 왜 교육과 직장에서는 완벽한 평등을 요구받으면서, 유독 번식에 있어서만은 비대칭적인 육체적 부담을 감당해야 하는가? 왜 동일한 속도로 달려야 하는 경쟁 트랙 위에서 나만 무거운 모래주머니(생물학적 부담)를 차야 하는가? 이 맹렬한 모순에 대한 역학적 반발과, 그것을 시스템적으로 보정해 달라는 거대한 에너지의 폭발이 바로 젊은 여성들의 '진보화'로 나타나는 것이다.
 
결국 놀랍게도, '젊은 세대의 젠더 양극화'와 '저출산'은 전혀 다른 두 개의 문제가 아니다. 하나는 정치학의 문제이고 하나는 인구학의 문제인 것처럼 보이지만, 디스턴스의 심해를 꿰뚫어 보면 이 둘은 "완벽한 대칭성(공정성)을 향해 질주하는 문명의 거대한 평활화 엔진이, 결코 깎여나가지 않는 번식의 생물학적 비대칭성이라는 단단한 암초와 정면충돌하며 박살 나고 있는 완벽히 동일한 '구조적 파열음'"일 뿐이다. 사회가 부유해질수록, 제도가 선진화될수록, 그리고 가장 치명적으로 '경쟁이 공정해질수록' 아이의 울음소리가 사라지고 남녀가 정치적으로 찢어지는 것은 결코 우연이 아니며 막을 수도 없는 열역학적 필연이다. 문명이 인위적 비대칭성을 매끄럽게 깎아내어 가장 완벽한 평활화를 달성하려 할수록, 그 깨끗한 유리창 위에 남은 마지막 생물학적 얼룩은 무견딜 수 없는 불공정으로 극대화되기 때문이다. 

인류는 과거의 낡은 계급과 불평등을 모두 갈아버리며 여기까지 왔다. 그러나 그 위대한 평활화 엔진의 칼날이 마침내 생명의 가장 근원적인 엇갈림(암수의 분화와 출산)이라는 뼈대에 도달했을 때, 문명은 더 이상 그 비대칭성을 깎아내지 못한 채 자신의 모멘텀을 잃고 요동치며 멈춰 서고 있는 것이다.
이것이 바로 공정성이 낳은 가장 잔혹한 우주적 역설이자, 저출산이 우리에게 던지는 궁극의 서늘한 진실이다.
